\chapter{Hola mundo}
\section{Configuracion}
Creamos una carpeta para el la practica de este libro y otra para react dentro del primero
\begin{lstlisting}[language=bash]
$ mkdir reactbook
$ mdir reactbook/react
\end{lstlisting}
A continuacion necesitamos agregar dos archivos uno para React el otro para el ReactDom
\begin{lstlisting}[language=bash]
$ curl -L https://unpkg.com/react@17/umd/react.development.js > ~/reactbook/
react/react.js
$ curl -L https://unpkg.com/react-dom@17/umd/react-dom.development.js > ~/react-
book/react/react-dom.js
\end{lstlisting}
Borra el @17 si ya contabas con una version de nodejs

\begin{lstlisting}[language=html]
<!doctype html>
<html>
  <head>
    <title>Hello React</title>
    <meta charset="utf-8" />
  </head>
  <body>
    <div id="app">
      <! -- my app renders here -->
      <h1>Hello World!</h1>
    </div>
    <script src="react/react.js"></script>
    <script src="react/react-dom.js"></script>
    <script>
      ReactDom.render(
        React.createElement("h1", null, "HelloWorld!"),
        document.getElementById("app"),
      );
    </script>
  </body>
</html>
\end{lstlisting}
Aqui podemos ver diferentes cosas
   
\begin{enumerate}
  \item El documento empieza siempre con <!Doctype html>
  \item El elemento de bloque padre seria <html></html> 
  \item Los otros dos elementos padres serian <head></head> y <body></body>
\end{enumerate}

a
